% Options for packages loaded elsewhere
\PassOptionsToPackage{unicode}{hyperref}
\PassOptionsToPackage{hyphens}{url}
\PassOptionsToPackage{dvipsnames,svgnames,x11names}{xcolor}
%
\documentclass[
]{article}
\usepackage{amsmath,amssymb}
\usepackage{iftex}
\ifPDFTeX
  \usepackage[T1]{fontenc}
  \usepackage[utf8]{inputenc}
  \usepackage{textcomp} % provide euro and other symbols
\else % if luatex or xetex
  \usepackage{unicode-math} % this also loads fontspec
  \defaultfontfeatures{Scale=MatchLowercase}
  \defaultfontfeatures[\rmfamily]{Ligatures=TeX,Scale=1}
\fi
\usepackage{lmodern}
\ifPDFTeX\else
  % xetex/luatex font selection
  \setmainfont[]{DejaVu Serif}
  \setmonofont[]{DejaVu Sans Mono}
\fi
% Use upquote if available, for straight quotes in verbatim environments
\IfFileExists{upquote.sty}{\usepackage{upquote}}{}
\IfFileExists{microtype.sty}{% use microtype if available
  \usepackage[]{microtype}
  \UseMicrotypeSet[protrusion]{basicmath} % disable protrusion for tt fonts
}{}
\makeatletter
\@ifundefined{KOMAClassName}{% if non-KOMA class
  \IfFileExists{parskip.sty}{%
    \usepackage{parskip}
  }{% else
    \setlength{\parindent}{0pt}
    \setlength{\parskip}{6pt plus 2pt minus 1pt}}
}{% if KOMA class
  \KOMAoptions{parskip=half}}
\makeatother
\usepackage{xcolor}
\usepackage[margin=1in]{geometry}
\usepackage{color}
\usepackage{fancyvrb}
\newcommand{\VerbBar}{|}
\newcommand{\VERB}{\Verb[commandchars=\\\{\}]}
\DefineVerbatimEnvironment{Highlighting}{Verbatim}{commandchars=\\\{\}}
% Add ',fontsize=\small' for more characters per line
\newenvironment{Shaded}{}{}
\newcommand{\AlertTok}[1]{\textcolor[rgb]{1.00,0.00,0.00}{\textbf{#1}}}
\newcommand{\AnnotationTok}[1]{\textcolor[rgb]{0.38,0.63,0.69}{\textbf{\textit{#1}}}}
\newcommand{\AttributeTok}[1]{\textcolor[rgb]{0.49,0.56,0.16}{#1}}
\newcommand{\BaseNTok}[1]{\textcolor[rgb]{0.25,0.63,0.44}{#1}}
\newcommand{\BuiltInTok}[1]{\textcolor[rgb]{0.00,0.50,0.00}{#1}}
\newcommand{\CharTok}[1]{\textcolor[rgb]{0.25,0.44,0.63}{#1}}
\newcommand{\CommentTok}[1]{\textcolor[rgb]{0.38,0.63,0.69}{\textit{#1}}}
\newcommand{\CommentVarTok}[1]{\textcolor[rgb]{0.38,0.63,0.69}{\textbf{\textit{#1}}}}
\newcommand{\ConstantTok}[1]{\textcolor[rgb]{0.53,0.00,0.00}{#1}}
\newcommand{\ControlFlowTok}[1]{\textcolor[rgb]{0.00,0.44,0.13}{\textbf{#1}}}
\newcommand{\DataTypeTok}[1]{\textcolor[rgb]{0.56,0.13,0.00}{#1}}
\newcommand{\DecValTok}[1]{\textcolor[rgb]{0.25,0.63,0.44}{#1}}
\newcommand{\DocumentationTok}[1]{\textcolor[rgb]{0.73,0.13,0.13}{\textit{#1}}}
\newcommand{\ErrorTok}[1]{\textcolor[rgb]{1.00,0.00,0.00}{\textbf{#1}}}
\newcommand{\ExtensionTok}[1]{#1}
\newcommand{\FloatTok}[1]{\textcolor[rgb]{0.25,0.63,0.44}{#1}}
\newcommand{\FunctionTok}[1]{\textcolor[rgb]{0.02,0.16,0.49}{#1}}
\newcommand{\ImportTok}[1]{\textcolor[rgb]{0.00,0.50,0.00}{\textbf{#1}}}
\newcommand{\InformationTok}[1]{\textcolor[rgb]{0.38,0.63,0.69}{\textbf{\textit{#1}}}}
\newcommand{\KeywordTok}[1]{\textcolor[rgb]{0.00,0.44,0.13}{\textbf{#1}}}
\newcommand{\NormalTok}[1]{#1}
\newcommand{\OperatorTok}[1]{\textcolor[rgb]{0.40,0.40,0.40}{#1}}
\newcommand{\OtherTok}[1]{\textcolor[rgb]{0.00,0.44,0.13}{#1}}
\newcommand{\PreprocessorTok}[1]{\textcolor[rgb]{0.74,0.48,0.00}{#1}}
\newcommand{\RegionMarkerTok}[1]{#1}
\newcommand{\SpecialCharTok}[1]{\textcolor[rgb]{0.25,0.44,0.63}{#1}}
\newcommand{\SpecialStringTok}[1]{\textcolor[rgb]{0.73,0.40,0.53}{#1}}
\newcommand{\StringTok}[1]{\textcolor[rgb]{0.25,0.44,0.63}{#1}}
\newcommand{\VariableTok}[1]{\textcolor[rgb]{0.10,0.09,0.49}{#1}}
\newcommand{\VerbatimStringTok}[1]{\textcolor[rgb]{0.25,0.44,0.63}{#1}}
\newcommand{\WarningTok}[1]{\textcolor[rgb]{0.38,0.63,0.69}{\textbf{\textit{#1}}}}
\setlength{\emergencystretch}{3em} % prevent overfull lines
\providecommand{\tightlist}{%
  \setlength{\itemsep}{0pt}\setlength{\parskip}{0pt}}
\setcounter{secnumdepth}{-\maxdimen} % remove section numbering
\usepackage{newunicodechar}
\newunicodechar{★}{\ensuremath{\star}}
\newunicodechar{✓}{\ensuremath{\checkmark}}
\newunicodechar{√}{\ensuremath{\surd}}
\newunicodechar{≈}{\ensuremath{\approx}}
\ifLuaTeX
  \usepackage{selnolig}  % disable illegal ligatures
\fi
\IfFileExists{bookmark.sty}{\usepackage{bookmark}}{\usepackage{hyperref}}
\IfFileExists{xurl.sty}{\usepackage{xurl}}{} % add URL line breaks if available
\urlstyle{same}
\hypersetup{
  colorlinks=true,
  linkcolor={Maroon},
  filecolor={Maroon},
  citecolor={Blue},
  urlcolor={Blue},
  pdfcreator={LaTeX via pandoc}}

\author{}
\date{}

\begin{document}

\hypertarget{chapter-3-why-sqrtd_k}{%
\section{\texorpdfstring{Chapter 3 --- Why
\(\sqrt{d_k}\)?}{Chapter 3 --- Why \textbackslash sqrt\{d\_k\}?}}\label{chapter-3-why-sqrtd_k}}

\begin{verbatim}
Type: theory + notebook
Ordering: theory-first
Depends on: 1 (scaled dot-product attention), 2 (softmax & temperature)
\end{verbatim}

\hypertarget{motivation}{%
\subsection{3.1 Motivation}\label{motivation}}

The scaling factor \(1/\sqrt{d_k}\) in scaled dot-product attention is
not cosmetic. It is the factor that keeps pre-softmax scores of order
one as the head dimension grows. Without it, softmax saturates at
initialization and gradients die. Why \(\sqrt{d_k}\) specifically ---
not \(d_k\), not \(\log d_k\), not a learned scalar --- is a variance
argument followed by a gradient argument.

\hypertarget{storyline}{%
\subsection{3.2 Storyline}\label{storyline}}

The mechanism is a three-link chain.

\begin{enumerate}
\def\labelenumi{\arabic{enumi}.}
\item
  \textbf{Variance in, variance out.} The inner product of two
  independent random vectors with i.i.d. unit-variance entries has
  variance equal to the dimension. Raw scores
  \(s_{ij} = \langle q_i, k_j \rangle\) therefore have standard
  deviation \(\sqrt{d_k}\), which grows unboundedly with head dimension.
  For \(d_k = 128\), scores live on the scale of \(\pm 11\).
\item
  \textbf{Softmax concentrates under large inputs.} From Chapter 2,
  softmax with a winner of margin \(\Delta\) gives
  \(p_{\text{win}} = e^\Delta / (e^\Delta + (n-1))\); once
  \(\Delta \gg \log(n-1)\), \(p_{\text{win}} \to 1\) and the
  distribution is effectively one-hot. Large-scale inputs push softmax
  into this saturated regime.
\item
  \textbf{Saturated softmax has zero Jacobian.}
  \(\partial p_i / \partial s_j = p_i(\delta_{ij} - p_j)\). When \(p\)
  is near one-hot every entry of this Jacobian is near zero. No gradient
  flows back to \(Q, K\), so attention cannot learn.
\end{enumerate}

Dividing by \(\sqrt{d_k}\) rescales the first link so pre-softmax inputs
have unit variance regardless of head dimension. Softmax then sees
inputs of moderate magnitude and the gradient survives.

\textbf{Regime notes.} The argument is sharpest \emph{at
initialization}. After training, learned \(W_Q, W_K\) restructure the
score distribution in ways the i.i.d. hypothesis does not capture (see
§3.3, Empirical observation 3.4). Measured pre-softmax score variance in
a trained model is typically below one even with the \(\sqrt{d_k}\)
factor present; the factor's job is to keep \emph{early} training
healthy, not to maintain unit variance throughout. The factor is a fixed
constant --- no train / val / inference distinction, no interaction with
dropout or masking.

\hypertarget{rigorous-core}{%
\subsection{3.3 Rigorous core}\label{rigorous-core}}

\textbf{Lemma 3.1 (Variance of an inner product).} Let
\(q, k \in \mathbb{R}^{d_k}\) have independent components \(q_i, k_i\)
with \(\mathbb{E}[q_i] = \mathbb{E}[k_i] = 0\) and
\(\text{Var}(q_i) = \text{Var}(k_i) = 1\), and \(q \perp k\) as random
vectors. Then
\[\mathbb{E}[\langle q, k \rangle] = 0, \qquad \text{Var}(\langle q, k \rangle) = d_k.\]
\emph{Sketch: mean by independence; variance by independence across
\(i\) and the factoring
\(\mathbb{E}[q_i^2 k_i^2] = \mathbb{E}[q_i^2]\mathbb{E}[k_i^2] = 1\).
Generalized in Exercise 3.1.}

\textbf{Proposition 3.2 (The \(\sqrt{d_k}\) correction).} Under the
hypotheses of Lemma 3.1,
\(\text{Var}(\langle q, k \rangle / \sqrt{d_k}) = 1\).

\textbf{Lemma 3.3 (Jacobian vanishing under saturation).} Let
\(p = \text{softmax}(s)\) for \(s \in \mathbb{R}^n\). Let
\(J_{ij} = \partial p_i / \partial s_j = p_i(\delta_{ij} - p_j)\). If
\(p \to e_k\) (one-hot at index \(k\)), then \(\|J\|_F \to 0\).
\emph{Sketch: case-split on \(i = k\) vs.~\(i \ne k\) --- every entry
vanishes individually.}

\textbf{Corollary 3.3.1.} Pre-softmax scores with standard deviation
\(\sqrt{d_k}\) place softmax in the saturating regime (Chapter 2, §2.4);
combined with Lemma 3.3, this means gradients at \(W_Q, W_K\) through an
unscaled attention layer vanish in expectation as \(d_k \to \infty\) at
initialization.

\textbf{Empirical observation 3.4 (The i.i.d. hypothesis breaks
post-training).} In a trained model the entries of \(q, k\) are not
i.i.d. --- learned \(W_Q, W_K\) concentrate mass in specific directions
dictated by the task. Pre-softmax score variance in trained attention is
typically \emph{below} one even with the \(\sqrt{d_k}\) factor present.
The scaling argument is sharp at init; after training it is
conservative. This motivates alternatives like QK-norm (Exercise 3.3)
that do not depend on the i.i.d. hypothesis.

\hypertarget{numerical-example-variance-scaling-is-a-distributional-fact}{%
\subsection{3.4 Numerical example --- variance scaling is a
distributional
fact}\label{numerical-example-variance-scaling-is-a-distributional-fact}}

A single hand computation of \(\langle q, k \rangle\) for fixed \(q, k\)
doesn't surface what Lemma 3.1 is saying --- the lemma is about the
\emph{distribution} of the inner product across i.i.d. samples, not any
particular value. The minimal experiment that reveals the mechanism is a
Monte-Carlo estimate of the variance across \(d_k\):

\begin{Shaded}
\begin{Highlighting}[]
\ImportTok{import}\NormalTok{ numpy }\ImportTok{as}\NormalTok{ np}

\NormalTok{M }\OperatorTok{=} \DecValTok{100\_000}
\ControlFlowTok{for}\NormalTok{ d\_k }\KeywordTok{in}\NormalTok{ [}\DecValTok{1}\NormalTok{, }\DecValTok{4}\NormalTok{, }\DecValTok{16}\NormalTok{, }\DecValTok{64}\NormalTok{, }\DecValTok{256}\NormalTok{, }\DecValTok{1024}\NormalTok{]:}
\NormalTok{    q }\OperatorTok{=}\NormalTok{ np.random.randn(M, d\_k)}
\NormalTok{    k }\OperatorTok{=}\NormalTok{ np.random.randn(M, d\_k)}
\NormalTok{    s }\OperatorTok{=}\NormalTok{ np.}\BuiltInTok{sum}\NormalTok{(q }\OperatorTok{*}\NormalTok{ k, axis}\OperatorTok{=}\DecValTok{1}\NormalTok{)}
    \BuiltInTok{print}\NormalTok{(}\SpecialStringTok{f"d\_k=}\SpecialCharTok{\{}\NormalTok{d\_k}\SpecialCharTok{:\textgreater{}5\}}\SpecialStringTok{  var(s) ≈ }\SpecialCharTok{\{}\NormalTok{s}\SpecialCharTok{.}\NormalTok{var()}\SpecialCharTok{:8.2f\}}\SpecialStringTok{   var(s/√d\_k) ≈ }\SpecialCharTok{\{}\NormalTok{(s }\OperatorTok{/}\NormalTok{ np.sqrt(d\_k))}\SpecialCharTok{.}\NormalTok{var()}\SpecialCharTok{:5.3f\}}\SpecialStringTok{"}\NormalTok{)}
\end{Highlighting}
\end{Shaded}

Expected: the first column tracks \(d_k\) linearly, the second stays
near \(1\). The notebook sweeps this across several random seeds and
measures the full distribution, not just the variance.

\hypertarget{mechanics-check}{%
\subsection{3.5 Mechanics check}\label{mechanics-check}}

\begin{itemize}
\tightlist
\item
  \textbf{Scale by \(1/\sqrt{d_k}\)} --- keeps pre-softmax score
  variance at \(O(1)\) regardless of head dimension (Lemma 3.1 +
  Proposition 3.2). Without this, softmax saturates at init and
  gradients die.
\item
  \textbf{Why \(\sqrt{d_k}\), not \(d_k\)} --- we want \emph{standard
  deviation} \(O(1)\), not variance \(O(1/d_k)\). Dividing by \(d_k\)
  shrinks scores too hard and collapses softmax toward uniform
  (temperature \(\to \infty\) in Chapter 2's sense), which is also
  uninformative.
\item
  \textbf{Why not a learned scalar at init} --- a learned scalar is
  initialized somewhere; starting far from \(1/\sqrt{d_k}\) creates the
  saturation problem this factor is designed to prevent. Learned scalars
  are a \emph{post-init} tool (see QK-norm, Exercise 3.3).
\item
  \textbf{Interaction with softmax, not with masking or dropout} --- the
  factor sits between \(QK^\top\) and softmax; it does not interact with
  the additive mask (which is applied to the scaled scores) or with
  attention dropout (which is applied to the softmax output). No train /
  val / inference asymmetry.
\end{itemize}

\hypertarget{exercises}{%
\subsection{3.6 Exercises}\label{exercises}}

\textbf{3.1 ★★} \emph{{[}Mechanism dissection --- general variances.{]}}
Let entries of \(q, k\) be i.i.d. with
\(\mathbb{E}[q_i] = \mathbb{E}[k_i] = 0\),
\(\text{Var}(q_i) = \sigma_q^2\), \(\text{Var}(k_i) = \sigma_k^2\).
Derive the scaling factor \(c(d_k, \sigma_q, \sigma_k)\) such that
\(\langle q, k \rangle / c\) has unit variance. What does this say about
the common convention of initializing both \(W_Q\) and \(W_K\) at the
same scale?

\textbf{3.2 ★★} \emph{{[}Failure-mode construction --- dependent \(q\)
and \(k\).{]}} Suppose \(q = W_Q x\) and \(k = W_K x\) for the
\emph{same} input \(x \in \mathbb{R}^d\), so \(q\) and \(k\) are not
independent. Let \(x\) have i.i.d. unit-variance entries. Derive
\(\text{Var}(\langle q, k \rangle)\) in terms of \(W_Q, W_K\). Does it
still scale like \(d_k\)? Construct explicit \(W_Q, W_K\) for which the
variance is \(\Theta(d_k^2)\) (much larger than \(d_k\)), and others for
which it is \(O(1)\) (much smaller). What structural feature of
\(W_Q^\top W_K\) controls the answer?

\textbf{3.3 ★★★} \emph{{[}Theorem about the numerics --- QK-norm.{]}}
QK-norm replaces \(QK^\top / \sqrt{d_k}\) with
\(\gamma \cdot \hat Q \hat K^\top\), where
\(\hat q_i = q_i / \|q_i\|_2\) row-wise (similarly \(\hat k_j\)), and
\(\gamma \in \mathbb{R}\) is learned.

\begin{enumerate}
\def\labelenumi{(\alph{enumi})}
\tightlist
\item
  Prove that \(\hat q^\top \hat k \in [-1, 1]\)
  \emph{deterministically}, independent of any distributional assumption
  on \(q, k\).
\item
  Conclude that pre-softmax scores are bounded by \(\pm \gamma\), and
  contrast this with the \(\sqrt{d_k}\) guarantee (which is an
  expectation, not a bound).
\item
  In Chapter 2's language, what quantity does \(\gamma\) play the role
  of? State the analogy precisely.
\item
  Give one situation in which QK-norm is strictly more robust than
  \(\sqrt{d_k}\)-scaling, and one in which it is strictly less
  expressive.
\end{enumerate}

\textbf{3.4 ★★★} \emph{{[}Predict-then-verify --- unscaled attention
across \(d_k\).{]}} The notebook will train a 1-layer transformer with
\emph{unscaled} attention on a simple copy task, sweeping
\(d_k \in \{4, 16, 64, 256\}\). For each \(d_k\), predict \emph{before
running} which of the following occurs, and justify from Lemmas 3.1 and
3.3 and Corollary 3.3.1:

\begin{enumerate}
\def\labelenumi{(\alph{enumi})}
\tightlist
\item
  training loss fails to decrease at all (plateau at the init-entropy
  floor);
\item
  training loss decreases slowly and converges to a worse optimum than
  the scaled baseline;
\item
  training loss decreases normally (the unscaled model learns fine);
\item
  training loss diverges (blows up) after some number of steps.
\end{enumerate}

Write your prediction as a \(4 \times 4\) table (rows: \(d_k\), columns:
outcomes (a)--(d), entries: your assigned probability summing to 1 per
row) with a one-line reason per row. State your overall confidence (low
/ medium / high, with reason).

\hypertarget{before-the-notebook-03_variance_scaling.ipynb}{%
\subsection{\texorpdfstring{3.7 Before the notebook
(\texttt{03\_variance\_scaling.ipynb})}{3.7 Before the notebook (03\_variance\_scaling.ipynb)}}\label{before-the-notebook-03_variance_scaling.ipynb}}

Write down the following predictions before opening the notebook. The
notebook's purpose is to test them.

\begin{enumerate}
\def\labelenumi{\arabic{enumi}.}
\tightlist
\item
  \textbf{Empirical variance of unscaled scores} as a function of
  \(d_k \in \{1, 4, 16, 64, 256, 1024\}\) --- shape, scaling exponent,
  numeric values at \(d_k = 1\) and \(d_k = 1024\).
\item
  \textbf{Empirical variance of scaled scores} --- shape, scaling
  exponent. State the expected deviation from 1 (concentration rate in
  \(M\), the number of Monte-Carlo samples).
\item
  \textbf{Softmax entropy} of a single row as a function of \(d_k\), for
  unscaled vs.~scaled inputs, with \(n = 32\) keys. State the limit of
  each as \(d_k \to \infty\) in one word.
\item
  \textbf{Max / mean softmax ratio} in a single row, as a function of
  \(d_k\), for unscaled vs.~scaled. Direction and limit.
\item
  \textbf{Frobenius norm of the softmax Jacobian} (evaluated at a random
  score vector) as a function of \(d_k\), unscaled vs.~scaled. Direction
  and the approximate rate of decay in the unscaled case (polynomial or
  exponential in \(d_k\)?).
\end{enumerate}

For each prediction: direction, order-of-magnitude estimate, confidence
level with a one-line reason. The notebook confirms or refutes each in
its Finding section.

\end{document}
